% ---------------------------------------------------------------------------
% Paper-format draft, version 2 (1 Oct 2026). Single file; compiles on
% Overleaf with pdflatex (run twice for references). Supersedes
% paper_draft.tex for the results sections: adds the pre-registered 7/2/4
% STGAT test and the market-making study.
%
% Provenance: every number comes from committed logs on branch
% fix/post-close-carry-forward (logs/v2_forecast/, runs/README.md,
% docs/RESEARCH_LOG.md, docs/planning/*preregistration*.md).
% "Exploratory" marks results that were not pre-registered.
% "Pilot" marks results from the May--Sep 2025 pilot window (paper_draft.tex).
%
% LaTeX notes carried over from data_window_memo.tex:
%   1. \textcent is text-mode only; use \cent (defined below) in maths.
%   2. Maths inside \textbf{} is not bolded; use \bm{} or \boldmath.
%   3. Check WARNINGS as well as errors when building.
% ---------------------------------------------------------------------------
\documentclass[11pt,a4paper]{article}

\usepackage[T1]{fontenc}
\usepackage{lmodern}
\usepackage[utf8]{inputenc}
\usepackage{textcomp}
\usepackage{amsmath}
\usepackage{amssymb}
\usepackage{bm}
\usepackage{booktabs}
\usepackage{tabularx}
\usepackage{multirow}
\usepackage[margin=2.4cm]{geometry}
\usepackage{enumitem}
\usepackage{xcolor}
\usepackage{url}
\usepackage{hyperref}
\hypersetup{colorlinks=true, linkcolor=black, citecolor=blue!50!black, urlcolor=blue!60!black}

\newcommand{\cent}{\text{\textcent}}
\newcommand{\x}{\ensuremath{\times}}
\newcolumntype{L}{>{\raggedright\arraybackslash}X}
\setlist[itemize]{itemsep=2pt, topsep=4pt}
\setlist[enumerate]{itemsep=2pt, topsep=4pt}
\setlength{\emergencystretch}{3em}

\title{\textbf{When Does a Learned Model Beat a No-Arbitrage Constraint?}\\[4pt]
\large Forecasting and Market Making in Kalshi Event Contracts}
\author{Justyn\\[2pt]
\small Final Year Project --- draft for supervisor review}
\date{1 October 2026}

\begin{document}
\maketitle

\begin{abstract}
Mutually exclusive, collectively exhaustive (MECE) baskets of event contracts must be priced to sum to \$1. This
no-arbitrage constraint gives free structure that a learned model could exploit. We ask where a spatio-temporal graph
attention network (STGAT) whose graph encodes the constraint adds value over the constraint itself, both as a
forecaster and as the basis of a trading strategy, using Kalshi trades and hourly top-of-book quotes from November
2024 to November 2025 with a pre-registered 7/2/4-month train/validation/test split.
(i) \emph{Baskets are priced coherently at quoted prices}: where every leg has a two-sided quote, the mid-price sum
is within $0.34\cent$ of \$1, and apparent dislocations in last-trade prices are mostly bid/ask print side.
(ii) \emph{The STGAT forecasts basket prices well}: on the held-out test months it has $0.635\x$ the error of
persistence and $0.922\x$ that of an optimal constraint-based shrinkage forecast. The margin comes from anticipating
settlement convergence; on unsettled baskets it does not beat a per-family constraint baseline.
(iii) \emph{Its forecasts do not translate into trading profit}: it cannot forecast next-hour quote moves, and it
does not improve a market maker whose profit is decided at settlement.
(iv) \emph{A learned market-making policy is profitable out of sample}: a gradient-boosted model that selects which
maker quotes to post earns $+3.55\%$ net of fees on the test months (95\% CI $+2.51\%$ to $+4.60\%$, 34{,}086
simulated fills) under a pessimistic fill model, while posting every quote loses money. The edge is not significant
under a strict model in which competing makers also improve their prices. Adding the STGAT's forecasts to the policy
reduces its return by 0.65 percentage points. We conclude that the learned graph model adds value in forecasting how
basket prices evolve, while the tradable edge lies in execution, which a model on simple market features captures.
\end{abstract}

% ===========================================================================
\section{Introduction}
% ===========================================================================

Prediction markets list contracts whose payoffs are tied together by construction. In a MECE basket, such as
``Which range will New York's high temperature fall in tomorrow?'', exactly one leg pays \$1, so the legs' prices
should sum to \$1. This no-arbitrage constraint is known exactly, before any data is seen.

That makes event contracts an unusual test bed for learned models. In most financial settings a model must discover
cross-asset structure from noisy data. Here the structure is given, and the research question becomes sharp:
\begin{quote}
\emph{Where does a learned model add value over the constraint itself, and where does the constraint suffice?}
\end{quote}

We study this on Kalshi, a CFTC-regulated exchange, with a spatio-temporal graph attention network (STGAT) whose
graph encodes the constraint: basket legs are linked through a basket hub. We evaluate it in two ways. First, as a
\emph{forecaster} of basket prices against baselines that use the constraint alone. Second, as the basis of a
\emph{trading strategy}, using a tape-driven market-making simulator with bounded fill assumptions and Kalshi's fee
schedule. Every evaluation reported as a test result was pre-registered and run once.

\paragraph{Contributions.}
\begin{enumerate}
  \item \textbf{A coherence measurement.} At quoted top-of-book prices, Kalshi baskets are coherent; the apparent
    dislocation in last-trade prices is microstructure (Section~\ref{sec:res-coherence}).
  \item \textbf{A pre-registered forecasting test.} The STGAT beats persistence and an optimal constraint-based
    shrinkage forecast on four held-out months; we locate the margin in settlement convergence
    (Section~\ref{sec:res-forecast}).
  \item \textbf{A diagnosis of why forecasting skill does not become trading profit.} A maker's profit is decided at
    settlement, a horizon at which no model beats the market price (Sections~\ref{sec:res-rule} and
    \ref{sec:res-settle}).
  \item \textbf{A learned market-making policy that is profitable out of sample} under pessimistic fills, with its
    limits stated, and a measurement of what the STGAT adds to it (Section~\ref{sec:res-policy}).
  \item \textbf{Negative results with mechanisms} for five further routes: cross-venue lead--lag, live weather
    observations, fill prediction, streaming data and a rule-based maker (Section~\ref{sec:res-other}).
\end{enumerate}

% ===========================================================================
\section{Background}
% ===========================================================================

\paragraph{Prediction markets.} Prices of binary event contracts are commonly read as probabilities
\cite{wolfers2004}. They are generally informative but show systematic biases, most prominently the
favourite--longshot bias \cite{snowberg2010}. On Kalshi, makers earn higher returns than takers, and high-priced
contracts win slightly more often than their price implies \cite{burgi2025}.

\paragraph{Basket dislocations and cross-venue pricing.} Saguillo et al.\ \cite{saguillo2025} estimate about \$40M
of realised arbitrage profit on Polymarket from April 2024 to April 2025, including MECE baskets priced away from
\$1; Polymarket was largely fee-free over that period. Across venues, Polymarket led price discovery for the 2024 US
presidential election \cite{aktug2026}. Kalshi's weather prices beat public numerical forecasts and lead them
\cite{crosier2026}, and frontier language-model agents lost money trading Kalshi in 2026 \cite{predictionarena2026}.

\paragraph{Market microstructure.} A trade executes at the bid or the ask, so trade prices ``bounce'' across the
spread even when the efficient price does not move \cite{roll1984}. Market makers earn the spread but lose to better
informed traders (adverse selection) \cite{glosten1985} and manage inventory risk \cite{avellaneda2008}. Replay
simulators for prediction markets must make assumptions about queue position and fees \cite{arora2026}.

\paragraph{Learning methods.} Graph attention \cite{velickovic2018} and spatio-temporal graph models
\cite{yu2018stgcn} combine cross-sectional structure with temporal dynamics; in microstructure, deep models usually
forecast \emph{changes} in a reference price rather than levels \cite{zhang2019deeplob}. For tabular data,
gradient-boosted trees \cite{ke2017lightgbm} remain competitive with, and often better than, deep networks
\cite{grinsztajn2022}.

% ===========================================================================
\section{Motivation}
% ===========================================================================

\subsection{From pairs trading to structural constraints}
The project began as a statistical pairs-trading study. A correlation- and half-life-based baseline on October 2025
data returned a Sharpe ratio of 5.52, but most selected pairs were economically unrelated, and the result is best
explained by multiple testing over tens of thousands of candidate pairs. Estimated relationships fit event contracts
poorly: most contracts live days or hours, binary payoffs have no stationary spread to revert to, and screening
millions of tickers guarantees spurious pairs. The scope therefore moved to relationships that need no estimation:
the MECE constraint, which is exact, known in advance and enforced at settlement.

\subsection{Why early positive results did not stand}
In the pilot window (May--September 2025), a model trained on last-trade prices first appeared to beat an optimal
constraint-based forecast by 15\% and to find tradeable dislocations. Both claims failed under scrutiny:
\begin{itemize}
  \item \textbf{Last-trade basket sums are dominated by microstructure.} 68.7\% of basket trade prints are ask-side;
    with a half-spread of $2.8\cent$ this predicts a $+6.3\cent$ excess on a 6-leg basket, against $+5.25\cent$
    measured. A model forecasting last-trade sums may be forecasting which side the next prints land on.
  \item \textbf{Pipeline leaks inflated the result:} resolved markets replayed after close, time-to-close from
    after-the-fact close times, basket membership chosen with hindsight, and baskets scored on a subset of legs.
\end{itemize}
We therefore rebuilt the pipeline causally, pre-registered every evaluation, scored against constraint-based
baselines rather than only persistence, and measured executable prices directly from quotes.

\subsection{From forecasting to trading}
A forecasting result alone does not show economic value. Taker strategies that trade basket dislocations were
unprofitable after fees at quoted prices in the pilot. Since makers, not takers, earn on Kalshi \cite{burgi2025},
we asked whether the STGAT can improve a \emph{market-making} strategy, and whether any model can make one
profitable out of sample.

% ===========================================================================
\section{Method}
% ===========================================================================

\subsection{Population and graph}
Contracts are nodes on a 2-hour snapshot grid. A \emph{MECE basket} is an event whose legs partition the outcome
space; membership is decided causally (a family qualifies only on instances known to have resolved at least 24h
earlier, and an event's own outcome is never used). Each basket has a hub node linked to its legs in both
directions. Node features summarise each 2-hour window (last price, VWAP, return, volatility, volume, order flow and
hub leg counts). Time-to-close is zeroed as an input.

\subsection{The STGAT}
The model has four parts: a type-specific input projection; a two-channel spatial attention layer (MECE hub--leg and
ladder); per-node \emph{causal} temporal self-attention over the snapshot sequence; and a residual forecast head
$\hat p_{t+h} = p_t + \Delta_h$ at horizons $h\in\{1,2,3,6,12\}$ snapshots (2--24 hours), with its final layer
zero-initialised so training starts exactly at persistence. An auxiliary head predicts the next-hour change in the
quoted mid. The loss combines per-node error with a basket-sum term; checkpoints are selected on validation MECE
error.

\subsection{Forecasting targets and baselines}
For basket $b$ with legs $L_b$, let $S_t=\sum_{i\in L_b} p_{i,t}$. A basket is scored only if every listed leg is
present; a leg that settles by $t+h$ is valued at its settlement (0 or 1). Baselines:
\begin{itemize}
  \item \textbf{Persistence:} $\hat S_{t+h}=S_t$.
  \item \textbf{Constraint oracle:} $\hat S_{t+h}=1+\alpha(S_t-1)$, with the single $\alpha$ that minimises squared
    error \emph{on the scored split}; deliberately optimistic.
  \item \textbf{Family oracle:} the same with one $\alpha$ per series \x\ horizon.
\end{itemize}

\subsection{Market-making simulator}
The simulator replays Kalshi's trade tape and hourly best bid/ask for MECE legs. At each hour it may post a bid
and/or an ask on a leg (10 contracts, improving the best price by $1\cent$ when the spread allows), cancels unfilled
orders after one hour, and holds every fill to settlement. Kalshi's maker fees are applied from the historical fee
schedule, with inventory limits of 50 contracts per leg and 200 per event. Whether a historical trade would have
filled our order depends on queue position, which a replay cannot observe, so every result is reported under three
fill models:
\begin{itemize}
  \item \textbf{Optimistic:} any trade at our price fills us.
  \item \textbf{Pessimistic} (the pre-registered verdict): a trade must go \emph{through} our price.
  \item \textbf{Strict:} competing makers also improve their price, so fills are rarer and more adverse.
\end{itemize}
Confidence intervals resample events (clusters).

\subsection{Strategies}
\begin{itemize}
  \item \textbf{Quote everything:} post both sides on every eligible leg (prices 1--99$\cent$, spread $\ge 2\cent$).
  \item \textbf{A10 (rules):} the most robust of 896 rule configurations searched on the development months:
    favourites only (order price 85--99$\cent$), spread $\ge 4\cent$, taker-only series, no quotes within 24h of
    close.
  \item \textbf{S4 / S5 (STGAT signal):} A10 with quotes skipped when the STGAT predicts an adverse move of more than
    half the spread (S4), or quoting only where the STGAT predicts a move (S5).
  \item \textbf{G (learned policy):} a gradient-boosted regressor predicts, for each candidate quote of ``quote
    everything'', its net P\&L per contract at settlement after fees; the quote is posted only if the prediction
    exceeds $1\cent$. The model is scikit-learn's \texttt{HistGradient\allowbreak{}Boosting\allowbreak{}Regressor}
    \cite{pedregosa2011} (400 iterations, learning rate 0.03, 31 leaves, at least 300 samples per leaf, L2 penalty
    1.0), with 14 features known at the quoting hour: order side and price, price relative to mid, spread, mid,
    trades in the last 24h, 24h mid change, basket size, legs quoted, basket sum of mids, best other leg's mid, hour,
    maker-fee status and category. Time-to-close is excluded because for sports markets the recorded close time is
    the actual match end.
  \item \textbf{G+S:} G plus six STGAT outputs at the quoting hour (the mid-move head and the forecast head's
    predicted change at each horizon, seed-averaged), from the latest snapshot no more than 2h old.
\end{itemize}

% ===========================================================================
\section{Experimental Setup}
% ===========================================================================

\begin{itemize}
  \item \textbf{Data.} Kalshi trades, market metadata and hourly quote candles, November 2024 to November 2025
    (November ends on the 25th). Earlier months have too few basket trades to be usable.
  \item \textbf{Split.} Strictly forward, 7/2/4 months: train November 2024--May 2025, validation June--July 2025,
    test August--November 2025.
  \item \textbf{STGAT tuning.} A pre-registered screen on validation MECE error selected the default configuration
    (embedding 32, 4 heads, learning rate $10^{-3}$, no dropout); four seeds were then trained. A variant with quote
    inputs was rejected by a pre-registered rule on validation.
  \item \textbf{Pre-registration.} Every test result below was specified in the repository before the test months
    were scored, and run once. The learned policies' models are hash-checked before scoring. Test months already
    used for the rule-based maker were reused for the learned policies; this is disclosed, and none of the learned
    policies' parameters used them.
  \item \textbf{Parlays.} From October 2025 Kalshi lists multi-game parlays (about 99{,}000 tickers that month).
    They belong to no basket; they were excluded from test scoring after a verified check that excluding edge-free
    nodes leaves MECE results byte-identical.
  \item \textbf{Leak checks for G.} With shuffled training labels the edge disappears; it survives removing
    time-to-close, fee status and category, and excluding sports series. An allow-all policy reproduces the
    no-policy simulation exactly.
\end{itemize}

% ===========================================================================
\section{Results}
% ===========================================================================

\subsection{Baskets are coherent at quoted prices (pilot)}\label{sec:res-coherence}

\begin{table}[htbp]
\centering
\caption{Basket sums on the pilot test month (September 2025): whole baskets with every leg live.}
\label{tab:coherence}
\begin{tabular}{lcc}
\toprule
 & all baskets ($n=7{,}532$) & every leg two-sided (69\%) \\
\midrule
Sum of last trades & $1.0405$ & $1.0261$ \\
Sum of quoted mids & $1.0185$ & $\bm{1.0034}$ \\
Sum of bids / asks & $0.944$ / $1.093$ & $0.967$ / $1.039$ \\
Selling at bids beats \$1 & 1.8\% of snapshots ($\approx 2\cent$) & --- \\
\bottomrule
\end{tabular}
\end{table}

Where every leg has a two-sided quote, the mid-price sum is within $0.34\cent$ of \$1 (Table~\ref{tab:coherence}).
In these baskets, 82--87\% of the last-trade excess is print side; the rest of the pooled excess sits in legs with no
bid. Executable dislocation is rare and small: selling at the bids beats \$1 in 1.8\% of snapshots by about
$2\cent$, before a taker fee of $1.3$--$5.8\cent$ per 6-leg basket.

\subsection{The STGAT forecasts basket prices}\label{sec:res-forecast}

\begin{table}[htbp]
\centering
\caption{Pre-registered test (August--November 2025): basket-sum forecast error relative to baselines. Four seeds,
mean (range); $n=109{,}642$ basket-horizon positions in $k=1{,}025$ events. Lower is better.}
\label{tab:forecast}
\begin{tabular}{lcc}
\toprule
Comparison & test & validation \\
\midrule
STGAT vs persistence & $\bm{0.635\x}$ ($0.567$--$0.683$) & $0.585\x$ \\
STGAT vs constraint oracle (single $\alpha=0.566$) & $\bm{0.922\x}$ ($0.824$--$0.991$) & $0.904\x$ \\
\quad seeds whose 95\% CI excludes 1 & 1 of 4 & --- \\
Unsettled baskets (80.1\%): vs single-$\alpha$ oracle & $0.991\x$ ($0.871$--$1.087$) & $0.916\x$ \\
Unsettled baskets: vs family oracle & $1.09\x$ ($0.959$--$1.196$) & $0.970\x$ \\
Next-hour mid-move head vs zero move & $1.03\x$ ($1.015$--$1.041$) & $1.007\x$ \\
Ladder pairs vs oracle (exploratory) & $0.991\x$ ($0.984$--$0.997$) & $0.994\x$ \\
\bottomrule
\end{tabular}
\end{table}

\begin{itemize}
  \item \textbf{The STGAT beats persistence clearly and the constraint oracle on average} (Table~\ref{tab:forecast}),
    but only one seed's CI excludes 1 against the oracle.
  \item \textbf{The margin comes from settlement convergence.} Baskets that settle within the horizon have a target
    of exactly \$1, which the model anticipates although time-to-close is not an input. On unsettled baskets it
    matches the single-$\alpha$ oracle and loses to the family oracle, the same pattern as in the pilot.
  \item \textbf{It cannot forecast next-hour quote moves}: every seed is worse than predicting no move.
\end{itemize}

\subsection{The STGAT cannot improve a rule-based maker}\label{sec:res-rule}

\begin{table}[htbp]
\centering
\caption{Pre-registered maker test (August--November 2025), pessimistic fills.}
\label{tab:rule}
\begin{tabular}{lcccc}
\toprule
Strategy & return & 95\% CI & net & fills / events \\
\midrule
A10 (rules) & $+1.80\%$ & $-48.5\%$ to $+62.4\%$ & $+\$1.87$ & 83 / 29 \\
A6 (rules, no 24h cut-off) & $-10.6\%$ & $-32.7\%$ to $+13.4\%$ & $-\$96.63$ & 1{,}852 / 754 \\
S4 (A10 + STGAT filter) & \multicolumn{4}{c}{identical to A10: the filter never fired} \\
S5 (STGAT-led) & $+24.2\%$ & $-36.7\%$ to $+90.7\%$ & $+\$14.94$ & 49 / 32 \\
\bottomrule
\end{tabular}
\end{table}

The rule-based maker A10 was positive on test but far from significant (Table~\ref{tab:rule}). Its development
return of $+21.7\%$ did not carry over, partly because its eligible opportunities collapsed from about 63 leg-days per
development month to 5--7 per month by October--November 2025. The STGAT did not change a single A10 trade. A
post-test diagnosis explains why: A10's profit is decided at settlement, not in the next hour. Its losing fills
earned $+\$62.59$ in their first hour on development data and lost $\$162.11$ by settlement. Skipping every fill with
an adverse first hour (perfect next-hour foresight) would have moved test P\&L only from $+\$1.87$ to $+\$3.07$,
whereas perfect foresight of the outcome would have given $+\$49.86$.

\subsection{Settlement outcomes are not forecastable beyond the price (exploratory)}\label{sec:res-settle}

If a maker's profit is decided at settlement, the useful forecast is the outcome itself. On all MECE legs in the
development months (validation: 9{,}790 leg-days, 891 events), no forecast beat the market price
(Table~\ref{tab:settle}). A gradient-boosted model with price, basket context and market features matched it; the
basket-normalised price (the constraint alone) was clearly worse, because the overround is not spread evenly across
legs.

\begin{table}[htbp]
\centering
\caption{Forecasting whether a leg settles YES (validation, exploratory). Lower is better.}
\label{tab:settle}
\begin{tabular}{lcc}
\toprule
Forecast & log-loss & Brier \\
\midrule
Market price (mid/100) & $\bm{0.4577}$ & $\bm{0.1498}$ \\
Price with isotonic calibration & $0.4603$ & $0.1507$ \\
Basket-normalised price, calibrated & $0.4758$ & $0.1579$ \\
Gradient-boosted model & $0.4589$ & $0.1500$ \\
\bottomrule
\end{tabular}
\end{table}

\subsection{A learned market-making policy is profitable out of sample}\label{sec:res-policy}

\begin{table}[htbp]
\centering
\caption{Pre-registered test (August--November 2025) of the learned policies, with validation (June--July 2025) in
brackets. Returns are net of fees on capital; CIs resample events.}
\label{tab:policy}
\small
\begin{tabularx}{\textwidth}{L c c c c c}
\toprule
Strategy & \shortstack{return,\\pessimistic} & 95\% CI & net & \shortstack{return,\\strict} & fills / events \\
\midrule
\textbf{G (learned policy)} & $\bm{+3.55\%}$ [$+3.73\%$] & $+2.51\%$ to $+4.60\%$ & $+\$3{,}748$ & $+0.39\%$ & 34{,}086 / 1{,}325 \\
G+S (with STGAT features) & $+2.90\%$ [$+3.38\%$] & $+1.88\%$ to $+3.85\%$ & $+\$3{,}235$ & $+0.01\%$ & 35{,}759 / 1{,}313 \\
Quote everything & $-0.34\%$ [$-0.36\%$] & $-0.77\%$ to $+0.10\%$ & $-\$1{,}133$ & $-3.07\%$ & 110{,}867 / 1{,}609 \\
A10 (rules) & $+1.80\%$ [$+105\%$] & $-48.5\%$ to $+62.4\%$ & $+\$1.87$ & $+28.5\%$ & 83 / 29 \\
\bottomrule
\end{tabularx}
\end{table}

\begin{itemize}
  \item \textbf{The learned policy passed its pre-registered verdict} (pessimistic return above 0 with the CI
    excluding 0): $+3.55\%$, $+\$3{,}748$ over 34{,}086 fills in 1{,}325 events (Table~\ref{tab:policy}). The same
    candidate quotes posted without the model lose money.
  \item \textbf{It held up out of sample}: $+3.73\%$ on validation and $+3.55\%$ on test, under both the optimistic
    ($+3.56\%$) and pessimistic models.
  \item \textbf{The edge is not significant under strict fills} ($+0.39\%$, CI $-0.96\%$ to $+1.78\%$): if competing
    makers also improve their prices, the policy's best fills are crowded out.
  \item \textbf{The STGAT does not help.} G+S minus G is $-0.65$ percentage points on test (CI $-1.40$ to $+0.09$)
    and $-0.35$ on validation, although STGAT features were available for 90.9\% of training fills. The STGAT's
    forecasts correlate $0.062$ with subsequent quote moves on its own training months but only $0.025$ on
    validation and $0.016$ on test: the policy learned to trust an in-sample-inflated signal whose out-of-sample
    information is near zero.
\end{itemize}

\subsection{Other routes}\label{sec:res-other}

\begin{table}[htbp]
\centering
\caption{Further routes to trading profit. ``Pre-reg.'' = pre-registered; otherwise exploratory, development data.}
\label{tab:other}
\small
\begin{tabularx}{\textwidth}{>{\raggedright\arraybackslash}p{5.2cm} L}
\toprule
Route & Outcome \\
\midrule
Kalshi $\rightarrow$ Polymarket catch-up (pre-reg., 2026 data, 768 matched soccer events)
  & Kalshi leads: {\boldmath $+0.70\cent$} per contract (CI $+0.47$ to $+0.97$, 4{,}320 trades), but Polymarket's sports
    taker fees turn it into $-0.53$ to $-1.35\cent$. On 2025 development data Polymarket did not lead Kalshi. \\
Weather brackets vs live station observations
  & Brackets made impossible by observed temperatures are repriced within the hour: at the first quote after,
    99.2\% bid below $2\cent$. \\
Predicting complete basket fills
  & Fills are predictable (AUC 0.91), but the predictable fills lose money: adverse selection. \\
Spotify \#1-song markets vs streaming data
  & Not testable: no free daily history for 2024--25. \\
\bottomrule
\end{tabularx}
\end{table}

Every route that relied on knowing more than the market failed (Table~\ref{tab:other}); the only lasting edge was in
execution.

% ===========================================================================
\section{Discussion and Limitations}
% ===========================================================================

\paragraph{Answer to the research question.} The learned graph model adds value in forecasting how basket prices
evolve, chiefly by anticipating settlement convergence, and the graph is what makes it competitive with the
constraint. It does not add value for trading. Kalshi prices are coherent at quoted prices, calibrated to outcomes,
and quick to absorb public information, so a model that forecasts prices finds nothing an execution strategy can use.
The profitable policy's edge is structural: which quotes are worth posting given the spread, the favourite--longshot
pattern and when informed traders are active. A gradient-boosted model on simple market features captures that,
consistent with evidence that boosted trees remain strong on tabular data \cite{grinsztajn2022}.

\paragraph{Limitations.}
\begin{itemize}
  \item \textbf{Fill assumptions.} Profitability holds under optimistic and pessimistic fills but not under strict
    competition; live queue position cannot be observed from the tape.
  \item \textbf{Backtest, not live trading.} Market impact, faster competitors and larger order sizes are not
    modelled; the simulator re-quotes hourly.
  \item \textbf{Sequential testing.} The learned policy was designed after the rule-based maker's test result was
    seen, and both were evaluated on the same four months. The policy's parameters never used those months, and its
    margin survives a correction for the six strategies tested, but a fresh forward period would make the test fully
    blind \cite{white2000, bailey2014}.
  \item \textbf{Scope.} The study covers November 2024 to November 2025, the period with enough basket trading to
    evaluate (earlier months have too few basket trades); this window was agreed in advance. Results describe one
    phase of a fast-growing market, concentrated in weather, crypto ranges and sports. The rapid growth from October
    2025 (for example, the arrival of parlays) suggests re-testing on later data before generalising, which we leave
    as possible future work.
  \item \textbf{In-sample STGAT features.} G+S used STGAT outputs that are in-sample on the training months, a caveat
    stated before the run; out-of-fold features could remove the harm but, given the out-of-sample correlations,
    are unlikely to add value.
\end{itemize}

\paragraph{Future work.}
\begin{itemize}
  \item \textbf{Forward test on later data.} Run the frozen learned policy, unchanged, on December 2025 onwards. This
    would give a fully blind test, addressing both the sequential-testing caveat and the scope of the window.
  \item \textbf{Model benchmark.} A pre-registered comparison on the same policy task: CatBoost, a modern deep tabular
    model, and the STGAT itself as the policy, each tuned on validation and run once on test.
\end{itemize}

% ===========================================================================
\section{Conclusion}
% ===========================================================================
A no-arbitrage constraint is a strong forecaster of basket prices, and a graph model built on it improves on it by
learning how baskets converge at settlement. That skill does not become trading profit, because the prices a trader
faces are already coherent and calibrated. Profit came instead from a learned execution policy, which is positive out
of sample under pessimistic fills but not under strict competition, and which the graph model's forecasts do not
improve.

% ===========================================================================
{\raggedright
\begin{thebibliography}{99}
% ===========================================================================
\bibitem{aktug2026} E.~Aktu\u{g} and O.~Torul. Price Discovery Across Political Prediction Markets: Evidence from the 2024 U.S. Presidential Election. Working paper, June 2026.
\bibitem{arora2026} A.~Arora and S.~Malpani. PredictionMarketBench: A SWE-bench-style framework for backtesting trading agents on prediction markets. arXiv:2602.00133, 2026.
\bibitem{avellaneda2008} M.~Avellaneda and S.~Stoikov. High-frequency trading in a limit order book. \emph{Quantitative Finance}, 8(3):217--224, 2008.
\bibitem{bailey2014} D.~H.~Bailey, J.~M.~Borwein, M.~L\'opez de Prado and Q.~J.~Zhu. Pseudo-mathematics and financial charlatanism: The effects of backtest overfitting on out-of-sample performance. \emph{Notices of the American Mathematical Society}, 61(5):458--471, 2014.
\bibitem{burgi2025} C.~B\"urgi, W.~Deng and K.~Whelan. Makers and Takers: The Economics of the Kalshi Prediction Market. CESifo / CEPR Discussion Paper DP20631, 2025.
\bibitem{crosier2026} A.~W.~Crosier. Prediction Markets Beat the Weather Forecast on Tomorrow's High Temperature. arXiv:2609.23969, 2026.
\bibitem{glosten1985} L.~R.~Glosten and P.~R.~Milgrom. Bid, ask and transaction prices in a specialist market with heterogeneously informed traders. \emph{Journal of Financial Economics}, 14(1):71--100, 1985.
\bibitem{grinsztajn2022} L.~Grinsztajn, E.~Oyallon and G.~Varoquaux. Why do tree-based models still outperform deep learning on typical tabular data? In \emph{NeurIPS Datasets and Benchmarks Track}, 2022.
\bibitem{ke2017lightgbm} G.~Ke, Q.~Meng, T.~Finley, T.~Wang, W.~Chen, W.~Ma, Q.~Ye and T.-Y.~Liu. LightGBM: A highly efficient gradient boosting decision tree. In \emph{NeurIPS}, 2017.
\bibitem{pedregosa2011} F.~Pedregosa et al. Scikit-learn: Machine learning in Python. \emph{Journal of Machine Learning Research}, 12:2825--2830, 2011.
\bibitem{predictionarena2026} Prediction Arena: Benchmarking AI models on real-world prediction markets. arXiv:2604.07355, 2026.
\bibitem{roll1984} R.~Roll. A simple implicit measure of the effective bid-ask spread in an efficient market. \emph{Journal of Finance}, 39(4):1127--1139, 1984.
\bibitem{saguillo2025} O.~Saguillo, V.~Ghafouri, L.~Kiffer and G.~Suarez-Tangil. Unravelling the Probabilistic Forest: Arbitrage in Prediction Markets. In \emph{Advances in Financial Technologies (AFT)}, LIPIcs, 2025. arXiv:2508.03474.
\bibitem{snowberg2010} E.~Snowberg and J.~Wolfers. Explaining the favorite-longshot bias: Is it risk-love or misperceptions? \emph{Journal of Political Economy}, 118(4):723--746, 2010.
\bibitem{velickovic2018} P.~Veli\v{c}kovi\'c, G.~Cucurull, A.~Casanova, A.~Romero, P.~Li\`o and Y.~Bengio. Graph attention networks. In \emph{ICLR}, 2018.
\bibitem{white2000} H.~White. A reality check for data snooping. \emph{Econometrica}, 68(5):1097--1126, 2000.
\bibitem{wolfers2004} J.~Wolfers and E.~Zitzewitz. Prediction markets. \emph{Journal of Economic Perspectives}, 18(2):107--126, 2004.
\bibitem{yu2018stgcn} B.~Yu, H.~Yin and Z.~Zhu. Spatio-temporal graph convolutional networks: A deep learning framework for traffic forecasting. In \emph{IJCAI}, 2018.
\bibitem{zhang2019deeplob} Z.~Zhang, S.~Zohren and S.~Roberts. DeepLOB: Deep convolutional neural networks for limit order books. \emph{IEEE Transactions on Signal Processing}, 67(11):3001--3012, 2019.
\end{thebibliography}
}

\end{document}
